\section{Hierarchical Skill Evolution}

\subsection{A hierarchy with different update clocks}

\method{} updates different abstraction levels at different rates. Low levels
change quickly because a mistake there touches only local evidence; high levels
change slowly because a mistake there propagates into every future task.
Table~\ref{tab:levels} lists the objects, their clocks, and the evidence each
promotion requires.

\begin{table}[t]
\caption{Skill-evolution objects and update clocks. The main experiments freeze
all training-derived transferable state before any downstream evaluation.}
\label{tab:levels}
\begin{center}
\small
\begin{tabular}{@{}c L{0.24\linewidth} L{0.24\linewidth} L{0.30\linewidth}@{}}
\toprule
Level & Object & Update clock & Evidence requirement \\
\midrule
L0 & Rollout trace & Every task attempt & Public trajectory and runtime outputs \\
L1 & Task evidence & Every task attempt & One task trace; verifier label (train only) \\
L2 & Repo skill candidate & Same-repo batch & Multiple evidence records, one repo \\
L3 & Behavior candidate & Cross-repo trigger & Repeated repo summaries across repos \\
L4 & Skill Writer policy & Slow policy window & Aggregated accept/reject histories \\
L5 & Skill Evaluator policy & Slow policy window & Verifier-calibrated false accepts/rejects \\
\bottomrule
\end{tabular}
\end{center}
\end{table}

Task evidence stores public fields only: task id, repo, problem summary,
reward, previous reward, case label, exception type, selected skills, touched and
edited paths, public test commands, failure signature, and diagnostic signals.
It feeds higher-level aggregation but is never exported as a reusable
\texttt{SKILL.md}.

\begin{figure}[t]
\centering
\begin{tikzpicture}[node distance=7mm and 12mm]
  \node[stage] (trace) {Agent\\trace};
  \node[stage, right=of trace] (evid) {Task\\evidence};
  \node[stage, right=of evid] (repo) {Repo\\candidate};
  \node[stage, right=of repo] (fail) {Behavior\\skill};
  % gates on the arrows
  \node[gate, above=4mm of $(evid)!0.5!(repo)$] (g1) {evaluator gate\\\scriptsize support\,/\,risk};
  \node[gate, above=4mm of $(repo)!0.5!(fail)$] (g2) {evaluator gate\\\scriptsize $\geq$\,3 repos\,/\,risk};
  \draw[flow] (trace) -- (evid);
  \draw[flow] (evid) -- node[below, font=\scriptsize, align=center]{same-repo\\batch} (repo);
  \draw[flow] (repo) -- node[below, font=\scriptsize, align=center]{cross-repo\\recurrence} (fail);
  \draw[flow, gray] (g1) -- ($(evid)!0.5!(repo)$);
  \draw[flow, gray] (g2) -- ($(repo)!0.5!(fail)$);
  \node[below=5mm of evid, font=\scriptsize, text width=2.4cm, align=center, gray]
    {public record,\\not deployed};
  \node[below=5mm of fail, font=\scriptsize, text width=2.6cm, align=center, gray]
    {primary\\transferable artifact};
\end{tikzpicture}
\caption{The promotion pipeline. A trace becomes task evidence, then a repo-level
candidate, then a cross-repo behavior skill. Task evidence is never deployed
directly; each upward step passes an evaluator gate that checks support, risk,
and validation. Repo promotion needs repeated same-repo evidence; behavior
promotion needs at least three distinct repositories.}
\label{fig:pipeline}
\end{figure}

\subsection{Writer and evaluator: generation versus gatekeeping}

The writer turns clusters of evidence into structured candidates. From a cluster
of same-repo evidence it emits a repo candidate with fields name, trigger,
evidence gate, actions, validation hint, abort condition, and support summary.
From a cross-repo cluster it emits a process-level behavior candidate with a
failure signature, trigger, actions, stop condition, and support summary:
\[
c_{\mathrm{repo}} = f_w(E_{\mathrm{repo}}), \qquad
c_{\mathrm{fail}} = f_w(E_{\mathrm{cross\text{-}repo}}).
\]
The structure is not decoration: it forces every candidate to expose its trigger,
support, actions, and risks so the evaluator can inspect them before anything is
deployed.

The evaluator is the gatekeeper that separates generation from deployment:
\[
\pi_e(c, E) \rightarrow (d, \hat r, \rho, z),
\]
with decision $d \in \{\textsc{accept}, \textsc{reject}, \textsc{revise},
\textsc{memory-only}\}$, proxy reward $\hat r$, confidence $\rho$, and risk flags
$z$. Hard safety risks (hidden verifier details, oracle or exact patches,
task-id memorization, private sandbox paths) force rejection regardless of
apparent utility. Repo candidates require repeated same-repo support; behavior
candidates require support from multiple distinct repositories. A single repo can
never originate a behavior skill.

We deliberately ship a \emph{deterministic} evaluator shell with explicit
thresholds and risk checks. This is an interface choice, not a limitation: it
makes the accept/reject contract auditable today, while exposing exactly the same
call signature that a learned evaluator, trained on SWEGym calibration
events, can later occupy without disturbing the rest of the loop.

\subsection{A fixed harness policy the learned policies cannot rewrite}

The learned writer and evaluator are explicitly forbidden from inventing their
own reward labels, update clocks, or promotion budgets. A fixed harness policy
maps verifier transitions and runtime diagnostics to case labels. A
\textbf{strong positive} is a previous attempt that failed and a current
attempt that passed; a \textbf{weak positive} is a current attempt that passed
without a paired improvement, or a previously passing attempt that stays
passing. A \textbf{strong negative} is a previous attempt that passed and a
current attempt that failed; a \textbf{weak negative} is a current attempt that
failed without a paired regression. A \textbf{diagnostic} covers missing reward,
timeout, setup error, verifier error, sandbox failure, or provider/runtime
failure. Cold-start solved traces are weak positives; unsolved traces are weak
negatives unless a diagnostic is present. Strong positives may add support to a
candidate but cannot by themselves create a global skill. Strong negatives
prioritize rollback, trigger narrowing, and evaluator false-accept penalties.
Diagnostics are kept out of semantic edit skills entirely and used only to
shape runtime or recover/validate safeguards.

\subsection{Training loop and the verifier boundary}

Training alternates rollout, evidence extraction, candidate generation, evaluator
gating, and validation (Figure~\ref{fig:verifier-boundary}, top lane). The
verifier supervises case labels, calibrates the evaluator, and gates candidate
states, and nothing more. At downstream evaluation (bottom lane) the writer,
evaluator, and transferable skills are frozen and only public trace signals enter
the update path.

\begin{figure}[t]
\centering
\begin{tikzpicture}[node distance=6mm and 10mm]
  % training lane
  \node[stage, fill=blue!4] (r1) {rollout};
  \node[stage, right=of r1] (e1) {task\\evidence};
  \node[stage, right=of e1] (c1) {repo/behavior\\candidates};
  \node[gate, right=of c1] (g1) {$\pi_e$ gate};
  \node[gate, right=of g1] (v1) {validation\\gate};
  \draw[flow] (r1)--(e1); \draw[flow] (e1)--(c1); \draw[flow] (c1)--(g1); \draw[flow] (g1)--(v1);
  \node[draw, fill=red!8, rounded corners, font=\scriptsize, above=5mm of c1] (ver)
    {SWEGym verifier $V$:\ labels, calibration, gates};
  \draw[flow, red!60!black, dashed] (ver) -- (c1);
  \draw[flow, red!60!black, dashed] (ver.east) to[out=0,in=90] (v1.north);
  \node[left=2mm of r1, font=\footnotesize\bfseries, rotate=90, anchor=south] {Train};

  % downstream lane
  \node[stage, fill=green!5, below=14mm of r1] (r2) {frozen\\$\pi_w,\pi_e$, skills};
  \node[stage, right=of r2] (a2) {agent\\rollout};
  \node[stage, right=of a2] (p2) {public\\evidence};
  \node[gate, right=of p2] (s2) {session\\skills only};
  \draw[flow] (r2)--(a2); \draw[flow] (a2)--(p2); \draw[flow] (p2)--(s2);
  \node[left=2mm of r2, font=\footnotesize\bfseries, rotate=90, anchor=south] {Test};
  % the boundary
  \draw[red!60!black, very thick, dashed] ($(r1.south west)!0.5!(r2.north west)+(-0.3,0)$)
     -- ($(v1.south east)!0.5!(s2.north east)+(0.3,0)$);
  \node[red!60!black, font=\scriptsize] at ($(s2.north east)+(0.1,0.9)$) {no $V$ below this line};
\end{tikzpicture}
\caption{The verifier boundary. \textbf{Top (train):} verifier $V$ supplies
outcome labels, evaluator calibration, and validation gates. \textbf{Bottom
(test):} the writer, evaluator, and transferable skills are frozen and only
public trace signals update memory. $V$ appears in the training lane and in final
reporting only---never on the downstream update path.}
\label{fig:verifier-boundary}
\end{figure}

\paragraph{Default update hyperparameters.}
Unless stated otherwise, repo updates use same-repo batches of 5 evidence events,
behavior updates trigger every 4 repo updates, a behavior promotion
requires at least 3 supporting repositories, and policy updates are capped at 3
bullet edits per policy window. Each repo batch emits at most one repo candidate;
each behavior trigger promotes at most two accepted behavior skills.
These values are the harness defaults and are treated as fixed protocol
hyperparameters, not tuned per benchmark.

\subsection{Validation gate}

Candidate skill states are scored on a repo-isolated SWEGym validation set via
\[
\mathrm{effective\_resolved\_rate}
= \frac{\#\,\text{resolved valid trials}}{\#\,\text{effective valid trials}},
\qquad
\mathrm{diagnostic\_rate}
= \frac{\#\,\text{diagnostic trials}}{\#\,\text{total trials}}.
\]
A candidate state is accepted only if it has enough effective validation trials,
stays under the diagnostic-rate threshold, and does not regress beyond a fixed
tolerance relative to the current accepted state. Otherwise the candidate version
is rejected or rolled back and the decision is logged. The gate exists precisely
to stop a plausible-looking candidate from silently trading pass rate for
regressions.
